\chapter{Estimation and Adaptive Control}
\label{chap:estimation}

\begin{learningobjective}
\begin{itemize}[leftmargin=*]
    \item Design observers and Kalman filters to reconstruct unmeasured states and reject sensor noise.
    \item Analyze innovation statistics to tune stochastic estimators.
    \item Explore model reference adaptive control (MRAC) through explicit update laws and simulation case studies.
\end{itemize}
\end{learningobjective}

\section{Concept Map for Beginners}
\begin{conceptroadmap}
\begin{enumerate}[leftmargin=*]
    \item Understand deterministic observers first—what they estimate, how fast they converge—and only then add noise to motivate Kalman filtering.
    \item Use innovation statistics as a feedback loop for your estimator, just as controller outputs are checked against performance metrics.
    \item Extend the estimation mindset to adaptive control: parameters become ``states'' we estimate online to keep performance robust.
\end{enumerate}
\end{conceptroadmap}

\section{Concept--Math--Engineering Loop}
\begin{table}[h]
    \centering
    \small
    \renewcommand{\arraystretch}{1.2}
    \begin{tabularx}{\linewidth}{|>{\raggedright\arraybackslash}p{0.28\linewidth}|>{\raggedright\arraybackslash}p{0.34\linewidth}|>{\raggedright\arraybackslash}X|}
        \hline
        \textbf{Concept} & \textbf{Mathematics} & \textbf{Engineering Practice} \\
        \hline
        Observer design & Pole placement of $A - L C$ & Use \texttt{place\_poles} and simulation hooks to confirm convergence with noisy measurements. \\
        \hline
        Kalman filtering & Riccati recursion, innovation covariance & Log $(t, r_k, S_k)$ from \texttt{EXPERIMENTS["kf\_lqr"]} to tune $Q$, $R$ empirically. \\
        \hline
        Adaptive control & Lyapunov analysis, MIT rule & Run \texttt{demo.py mrac} and tutorial plots to visualise parameter learning against design goals. \\
        \hline
    \end{tabularx}
    \caption{Estimation and adaptation concepts tied to their mathematical tools and executable experiments.}
\end{table}

\section{Deterministic State Observers}
\begin{beginnerguide}
    \item[\textbf{What?}] Luenberger observers that reconstruct full state vectors using known dynamics and measurements.
    \item[\textbf{Why?}] Many systems lack direct state sensing; observers provide the feedback signals controllers require.
    \item[\textbf{When?}] When sensors measure only a subset of states but the system is observable.
    \item[\textbf{How?}] Select observer poles, compute gain $L$, and simulate error dynamics to verify convergence speed.
    \item[\textbf{Example}] Place observer poles five times faster than closed-loop poles using \texttt{place\_poles} and inspect error decay in simulation.
\end{beginnerguide}
\subsection{Derivation}
For partially measured systems, the discrete Luenberger observer
\[
    \hat{x}_{k+1} = A \hat{x}_k + B u_k + L (y_k - C \hat{x}_k)
\]
reconstructs the state. The error $e_k = x_k - \hat{x}_k$ evolves as $e_{k+1} = (A - L C) e_k$. Choose $L$ such that the eigenvalues of $A - L C$ lie well inside the unit circle (typically 2--5 times faster than the closed-loop poles).

\subsection{Design Steps}
\begin{enumerate}[leftmargin=*]
    \item Specify desired observer poles $\Lambda_{\text{obs}}$.
    \item Form the dual system $(A^\top, C^\top)$ and compute $L^\top$ via pole placement (e.g., \texttt{scipy.signal.place\_poles}).
    \item Simulate the observer with noisy measurements to ensure convergence speed and noise rejection align with expectations.
\end{enumerate}

\subsection{Code Mapping}
\begin{table}[h]
    \centering
    \small
    \renewcommand{\arraystretch}{1.2}
    \begin{tabularx}{\linewidth}{|>{\raggedright\arraybackslash}p{4cm}|X|X|}
        \hline
        \textbf{Python} & \textbf{Formula} & \textbf{Explanation} \\
        \hline
        \texttt{scipy.signal.place\_poles}\\\texttt{(A.T, C.T, poles).gain\_matrix.T} & $L$ & Generates the observer gain from desired poles. \\
        \hline
        \texttt{simulate\_discrete\_system} with observer-in-the-loop & $e_{k+1} = (A - L C) e_k$ & Validates convergence numerically. \\
        \hline
    \end{tabularx}
    \caption{Luenberger observer implementation.}
\end{table}

\section{Kalman Filtering}
\begin{beginnerguide}
    \item[\textbf{What?}] Probabilistic state estimation that minimises mean-square error under Gaussian noise assumptions.
    \item[\textbf{Why?}] It weights model predictions and measurements optimally, delivering robust estimates and confidence bounds.
    \item[\textbf{When?}] When process and measurement noise are significant or uncertainty quantification is needed.
    \item[\textbf{How?}] Propagate $(\hat{x},P)$ through prediction/correction steps and tune $Q$, $R$ using innovation statistics.
    \item[\textbf{Example}] Run \texttt{EXPERIMENTS["kf\_lqr"].runner(return\_logs=True)} to log innovations and adjust $Q$ until the normalised innovation squared averages near the measurement dimension.
\end{beginnerguide}
\subsection{Stochastic Model}
Assume
\begin{align}
    x_{k+1} &= A x_k + B u_k + w_k, & w_k &\sim \mathcal{N}(0, Q), \\
    y_k &= C x_k + v_k, & v_k &\sim \mathcal{N}(0, R),
\end{align}
with mutually independent noises. The Kalman recursion is
\begin{align*}
    \hat{x}_{k|k-1} &= A \hat{x}_{k-1|k-1} + B u_{k-1},\\
    P_{k|k-1} &= A P_{k-1|k-1} A^\top + Q,\\
    K_k &= P_{k|k-1} C^\top (C P_{k|k-1} C^\top + R)^{-1},\\
    \hat{x}_{k|k} &= \hat{x}_{k|k-1} + K_k (y_k - C \hat{x}_{k|k-1}),\\
    P_{k|k} &= (I - K_k C) P_{k|k-1}.
\end{align*}
The innovation $r_k = y_k - C \hat{x}_{k|k-1}$ has covariance $S_k = C P_{k|k-1} C^\top + R$ and serves as a diagnostic statistic.

\subsection{Tuning Guidelines}
\begin{itemize}[leftmargin=*]
    \item Increase $Q$ to trust the model less and rely more on measurements (useful when unmodeled disturbances dominate).
    \item Increase $R$ to trust the measurements less (useful when sensors are noisy).
    \item Monitor the normalized innovation squared (NIS) statistic $\tilde{y}_k^T S_k^{-1} \tilde{y}_k$; its average should approach the measurement dimension.
\end{itemize}

\subsection{Code Mapping and Logging}
\begin{table}[h]
    \centering
    \small
    \renewcommand{\arraystretch}{1.2}
    \begin{tabularx}{\linewidth}{|>{\raggedright\arraybackslash}p{4cm}|X|X|}
        \hline
        \textbf{Python} & \textbf{Formula} & \textbf{Explanation} \\
        \hline
        \texttt{predict()} & $\hat{x}_{k|k-1}$, $P_{k|k-1}$ & Method of the filter object; propagates the state and covariance. \\
        \hline
        \texttt{update()} & $K_k$, $r_k$, $P_{k|k}$ & Applies the correction step and returns the innovation together with the updated covariance. \\
        \hline
        \texttt{runner()} & $(t, r_k, S_k)$ & Invoke on \texttt{EXPERIMENTS["kf\_lqr"]} with \texttt{return\_logs=True} to capture the innovation statistics used to tune $Q$, $R$. \\
        \hline
    \end{tabularx}
    \caption{Kalman filter code correspondence.}
\end{table}

\subsection{Case Study}
Select $Q = \operatorname{diag}(10^{-4}, 5\times 10^{-5})$, $R = 4\times 10^{-4}$ for the mass--spring--damper with position measurements only. Running \texttt{python demo.py kf\_lqr --no-plot} yields steady-state gains and innovation statistics. Adjust $Q$ by an order of magnitude to see the trade-off between responsiveness and noise amplification; inspect the logged $S_k$ to verify tuning.

\section{Model Reference Adaptive Control}
\begin{beginnerguide}
    \item[\textbf{What?}] Adaptive laws that adjust controller parameters online so the plant follows a reference model.
    \item[\textbf{Why?}] MRAC counters parameter drift and unmodelled dynamics without manual retuning.
    \item[\textbf{When?}] After ensuring persistent excitation and when plant parameters are uncertain or changing.
    \item[\textbf{How?}] Define the reference model, choose an adaptation law (e.g., MIT rule), and monitor Lyapunov-based error metrics.
    \item[\textbf{Example}] Execute \texttt{python demo.py mrac} to observe $\theta_k$ convergence and tracking error shrinking toward the reference model output.
\end{beginnerguide}
\subsection{Derivation}
MRAC enforces $y_k \rightarrow y_{m,k}$ where the reference model satisfies
\begin{equation}
    y_{m,k+1} = a_m y_{m,k} + b_m r_k.
\end{equation}
For the first-order plant $y_{k+1} = a y_k + b u_k$, use the control law
\begin{equation}
    u_k = \theta_1 y_k + \theta_2 r_k
\end{equation}
with parameter update via the MIT rule
\begin{equation}
    \theta_{k+1} = \theta_k - \gamma \phi_k e_k,
\end{equation}
where $\phi_k = [y_k\; r_k]^\top$ and $e_k = y_k - y_{m,k}$. Define $J_k = \tfrac{1}{2} e_k^2$; the gradient step
\[
    \theta_{k+1} = \theta_k - \gamma \frac{\partial J_k}{\partial \theta}
                 = \theta_k - \gamma \left(\frac{\partial e_k}{\partial \theta}\right) e_k
\]
reduces to the MIT rule because $\partial e_k / \partial \theta = \phi_k$.

\subsection{Stability Insight}
Define a Lyapunov function candidate
\begin{equation}
    V = \frac{1}{2} e_k^2 + \frac{1}{2 \gamma} (\theta_k - \theta^*)^T (\theta_k - \theta^*),
\end{equation}
where $\theta^*$ denotes ideal parameters that achieve perfect tracking. Under standard MRAC assumptions (matching conditions, persistent excitation), the adaptation law ensures $V$ decreases, leading to bounded tracking error.

\subsection{Code Mapping}
\begin{table}[h]
    \centering
    \small
    \renewcommand{\arraystretch}{1.2}
    \begin{tabularx}{\linewidth}{|>{\raggedright\arraybackslash}p{4cm}|X|X|}
        \hline
        \textbf{Python} & \textbf{Formula} & \textbf{Explanation} \\
        \hline
        \texttt{MRACController.step()} & $u_k = \theta^\top \phi_k$ & Applies the adaptive control law. \\
        \hline
        \texttt{MRACController.theta} & $\theta_k$ & Stores parameter estimates; log to analyse convergence. \\
        \hline
    \end{tabularx}
    \caption{MRAC algorithm in \texttt{scripts/controllers.py}.}
\end{table}

\subsection{Simulation Practice}
Run \texttt{python demo.py mrac}. The experiment discretizes a first-order plant with $a=-0.4$, $b=1$, uses $a_m=0.97$, $b_m=0.03$, and $\gamma=0.2$. Record steady-state error, overshoot, and parameter evolution. Adjust $\gamma$ to explore convergence versus noise sensitivity; add parameter logging to visualise $\theta_k$ settling.

\section{Tutorial: Adaptive Control Walkthrough}
\begin{beginnerguide}
    \item[\textbf{What?}] A scripted tour through adaptive tracking and parameter convergence figures.
    \item[\textbf{Why?}] Visualising trajectories helps interpret adaptation speed and robustness to noise.
    \item[\textbf{When?}] After learning the equations, as a sandbox for experimenting with adaptation gains.
    \item[\textbf{How?}] Run \tutorialcmd[--output-dir figures]{adaptive} and inspect tracking and parameter plots while varying $\gamma$.
    \item[\textbf{Example}] Increase $\gamma$ to 0.4, rerun the tutorial, and note faster convergence but amplified parameter oscillations.
\end{beginnerguide}
\begin{enumerate}[leftmargin=*]
    \item Produce the adaptive-control figures with\\
    \tutorialcmd[--output-dir figures]{adaptive}.
    \item Inspect Figure~\ref{fig:adaptive-tracking} to see how the plant output matches the reference model; relate the transient to the Lyapunov design insights in \textcite{krstic1995nonlinear}.
    \item Examine Figure~\ref{fig:adaptive-parameters} to monitor the evolution of $\theta_1$ and $\theta_2$. Adjust the adaptation gain $\gamma$ inside \texttt{tutorial\_adaptive} to visualise the trade-off between convergence speed and robustness.
\end{enumerate}
Appendix~\ref{chap:tooling} lists the supported \texttt{demo.py} experiments and plotting options you can reuse while exploring these scenarios.
\IfFileExists{../figures/adaptive_tracking.png}{
    \begin{figure}[h]
        \centering
        \includegraphics[width=0.75\linewidth]{../figures/adaptive_tracking.png}
        \caption{MRAC tracking performance generated by \tutorialcmd{adaptive}.}
        \label{fig:adaptive-tracking}
    \end{figure}
}{
    \begin{figure}[h]
        \centering
        \fbox{\parbox{0.7\linewidth}{Run \tutorialcmd{adaptive} to create the MRAC tracking plot, then recompile.}}
        \caption{Placeholder for the MRAC tracking plot.}
        \label{fig:adaptive-tracking}
    \end{figure}
}
\IfFileExists{../figures/adaptive_parameters.png}{
    \begin{figure}[h]
        \centering
        \includegraphics[width=0.7\linewidth]{../figures/adaptive_parameters.png}
        \caption{Parameter adaptation trajectories illustrating online learning.}
        \label{fig:adaptive-parameters}
    \end{figure}
}{
    \begin{figure}[h]
        \centering
        \fbox{\parbox{0.7\linewidth}{Generate the parameter trajectories with \tutorialcmd{adaptive}, then recompile.}}
        \caption{Placeholder for the MRAC parameter plot.}
        \label{fig:adaptive-parameters}
    \end{figure}
}

\section{Handling Missing or Delayed Measurements}
\begin{beginnerguide}
    \item[\textbf{What?}] Strategies for keeping estimators stable under packet drops, outliers, or sensor latency.
    \item[\textbf{Why?}] Real-world systems rarely deliver perfect data; robustness to measurement issues is essential.
    \item[\textbf{When?}] During estimator deployment, especially over networks or with low-cost sensors.
    \item[\textbf{How?}] Propagate without updates during dropouts, inflate covariance, and gate outliers.
    \item[\textbf{Example}] Modify \texttt{simulate\_discrete\_system} to drop every fifth measurement and observe how covariance inflation maintains state estimate accuracy.
\end{beginnerguide}
Modify \texttt{simulate\_discrete\_system} to include packet drops or sensor delays. Strategies to maintain estimator performance include:
\begin{itemize}[leftmargin=*]
    \item \textbf{Prediction during dropouts}: propagate the estimate using the process model when measurements are unavailable.
    \item \textbf{Covariance inflation}: increase $P$ during prolonged dropouts to reflect growing uncertainty.
    \item \textbf{Outlier rejection}: ignore innovations exceeding a specified threshold to handle bad data.
\end{itemize}

\section{Hands-On Practice}
\begin{beginnerguide}
    \item[\textbf{What?}] Exercises that apply observers, Kalman filters, and adaptive controllers to scripted experiments.
    \item[\textbf{Why?}] Practice reveals tuning nuances and builds confidence before tackling custom plants.
    \item[\textbf{When?}] After studying the theory—use the tasks as immediate reinforcement.
    \item[\textbf{How?}] Run the experiment harness, log results, adjust gains/noise levels, and compare against analytical expectations.
    \item[\textbf{Example}] Toggle measurement dropouts while plotting innovations to see how covariance inflation affects the Kalman filter’s residual statistics.
\end{beginnerguide}
\begin{enumerate}[leftmargin=*]
    \item Design a Luenberger observer using desired poles; verify $A - L C$ eigenvalues numerically.
    \item Tune the Kalman filter by logging innovations:
\begin{lstlisting}[language=python]
result = EXPERIMENTS["kf_lqr"].runner(return_logs=True)
innovation_log = result.metadata["innovation_log"]
\end{lstlisting}
    Adjust $Q$, $R$ until the residual statistics match the theoretical covariance.
    \item Use the MRAC experiment to visualise parameter convergence for different $\gamma$ values; relate overshoot changes to the Lyapunov analysis.
    \item Introduce measurement dropouts and confirm covariance inflation keeps the estimator responsive.
\end{enumerate}

\begin{takeaway}
\begin{itemize}[leftmargin=*]
    \item Observers and Kalman filters reconstruct hidden states and mitigate sensor noise; tune them using innovation statistics.
    \item Adaptive control compensates for parameter drift but must be monitored to avoid instability.
    \item Combined estimation and adaptation enable robust performance in the face of uncertainty and measurement imperfections.
\end{itemize}
\end{takeaway}

\section{Module Review}
\begin{itemize}[leftmargin=*]
    \item \textbf{Observer essentials}: choose observer poles (or gains) at least twice the closed-loop bandwidth, use Lyapunov criteria to confirm convergence.
    \item \textbf{Kalman tuning}: base $Q$, $R$ on physical insight and refine via innovation statistics.
    \item \textbf{Adaptive laws}: MIT update $\theta_{k+1} = \theta_k - \gamma \phi_k e_k$ balances tracking accuracy and robustness; monitor parameter trajectories.
    \item \textbf{Practice}: run \tutorialcmd{adaptive} to explore noise, dropout, and gain effects.
    \item \textbf{Reading}: \textcite{krstic1995nonlinear} for rigorous adaptive proofs and \textcite{ljung2012control} for identification-integrated estimation.
\end{itemize}

\section{Keyword Review}
\begin{vocabulary}
\begin{itemize}[leftmargin=*]
    \item \textbf{Luenberger observer}: deterministic estimator with error dynamics $A - L C$ tuned via pole placement.
    \item \textbf{Innovation}: difference $r_k = y_k - C \hat{x}_{k|k-1}$ summarising new measurement information.
    \item \textbf{Kalman gain}: optimal weighting $K_k = P_{k|k-1} C^\top (C P_{k|k-1} C^\top + R)^{-1}$ balancing model and measurements.
    \item \textbf{Normalized innovation squared (NIS)}: statistic $r_k^\top S_k^{-1} r_k$ used to diagnose filter tuning.
    \item \textbf{Model reference adaptive control (MRAC)}: scheme forcing plant output toward a reference model via online parameter updates.
    \item \textbf{MIT rule}: gradient-descent update $\theta_{k+1} = \theta_k - \gamma \phi_k e_k$ minimizing tracking error.
\end{itemize}
\end{vocabulary}

For motivated readers, \textcite{krstic1995nonlinear} expand these ideas to nonlinear plants, while \textcite{ljung2012control} show how estimation theory integrates with system identification workflows.
